\begin{proposition}[The global height formula]\hypertarget{tag-000V}{}\label{ari:C.6.1}
\StatementRecord{ari:C.6.1}{proved}{000V}{\SeriesRef{ari:C.7.1}; \citet[Section 2.1]{demarco2017}}

For every number field $K$, every $t\in K\setminus\{0,1\}$,
and the section just specified,
\begin{equation}\hypertarget{tag-0039}{}\label{ari:eq:ch-global}
 \sum_{\sigma:K\hookrightarrow\C}\log N(P_{\sigma(t)})
 =\sum_{v\nmid\infty}n_v\lambda_v(P_t)
          -[K:\Q]\widehat h_O(P_t).
\end{equation}
Here the local heights are invariant under change of Weierstrass
model and satisfy, when $nP\ne O$,
\begin{equation}\hypertarget{tag-003B}{}\label{ari:eq:ch-local-mult}
 \lambda_v(nP)=n^2\lambda_v(P)-\log|\psi_n(P)|_v
                   +\frac{n^2-1}{12}\log|\Delta|_v.
\end{equation}
At a complex embedding $\lambda_v(P)=-\log N(P)$; complex
embeddings are counted separately in \eqref{ari:eq:ch-global}.
For rational $t$, put $A_t=\exp(\sum_p\lambda_p(P_t))$.
Then $A_t$ is a positive real algebraic number and
\[
 N(P_t)=A_t e^{-\widehat h_O(P_t)},\qquad
 D_{\rm part}(t)=\frac6\pi
       \left(\log\frac{A_t}{2\sqrt3}-\widehat h_O(P_t)\right).
\]
Thus $N(P_t)$ is algebraic if and only if
$e^{\widehat h_O(P_t)}$ is algebraic.
\end{proposition}


\begin{proof}
\begin{enumerate}[label=\arabic*.,leftmargin=*]
\item The sigma expression for the archimedean local height is
\[
 \lambda_v(P)=-\log|\sigma(u)|
       +\frac12\Re(u\eta(u))-\frac1{12}\log|\Delta|.
\]
The sigma--theta relation and the Legendre relation identify this
with $-\log N$ exactly, including its constant; this is also the
normalization in \cref{ari:C.4.1}. The division identity
$\sigma(nu)=(-1)^{n-1}\psi_n(P)\sigma(u)^{n^2}$ gives
\eqref{ari:eq:ch-local-mult}. At finite places the same identity is
the defining multiplication law of the canonical N\'eron function.
Its normalization is uniquely determined by continuity away from
$O$, its logarithmic pole at $O$, and the multiplication law.
These are the three properties printed in
\citet[Section~2.1, arXiv v2, pp.~5--6]{demarco2017}; their duplication law includes
the discriminant term $\frac14\log|\Delta|_v$.
For the homothety
$(x,y,\Delta,\psi_n)\mapsto
(c^{-2}x,c^{-3}y,c^{-12}\Delta,c^{1-n^2}\psi_n)$,
the change in the two logarithmic terms of
\eqref{ari:eq:ch-local-mult} is
$-(1-n^2)\log|c|_v-(n^2-1)\log|c|_v=0$.
The pole normalization is transported by the same formal-parameter map.

\item Sum the duplication law over places. The product formula cancels
the division polynomial and discriminant terms, so that the sum
$H(P)=[K:\Q]^{-1}\sum_vn_v\lambda_v(P)$ satisfies
$H(2P)=4H(P)$. The pole normalization implies
$H(P)=\frac12h(x(P))+O_E(1)$. Iterating multiplication by two
therefore identifies $H$ with $\widehat h_O$. Moving the
archimedean terms to the other side gives \eqref{ari:eq:ch-global}.
For a finite place, a multiple of $P$ lies in the identity component
of a minimal model. There the local height is a rational multiple
of $\log p$; the multiplication law shows the same for $P$.
Only finitely many places contribute. This proves the assertion
about $A_t$ and both of its consequences.
\end{enumerate}
\end{proof}